%!TEX root = paper.tex

\section{Properties of Lewis weights - proofs}
\label{sec:properties-proofs}
This appendix contains proofs of the properties stated in Section~\ref{sec:properties}.

\begin{lemma}
\label{lem:equivstable}
Parts \ref{def:stable1} and \ref{def:stable2} of Definition~\ref{def:stable} are equivalent.
\end{lemma}
\begin{proof}
Given $\alpha$-almost Lewis weights $\ww$, define the errors $\ee$ as
\begin{equation*}
\ee_i = \aa_i^T \left( \AA^T \WW^{1 - 2/p} \AA \right)^{-1} \aa_i \ww_i^{-2/p}.
\end{equation*}
Then define $\AA'$ as $\EE^{1/p-1/2} \AA$ and $\BB'$ as $\WW^{1/2-1/p} \AA$. $\AA'$ has
Lewis weights $\ee_i \ww_i$, which reweight it to $\BB'$.  $\AA$,
when reweighted by the true weights $\wwbar$, will give $\BB = \WWbar^{1-2/p} \AA$.
$\AA$ is within a factor of $\alpha^{|1/2-1/p|}$ of $\AA'$; part \ref{def:stable1} then would imply
that $\BB$ is within a $\alpha^{c |1/2-1/p|}$ factor of $\BB$ and equivalently that
\begin{equation*}
\wwbar_i \approx_{\alpha^c} \ww_i
\end{equation*}
giving \ref{def:stable2}.
This argument works in reverse as well (interpreting the Lewis weights for any $\AA'$ as scaled
approximate Lewis weights for $\AA$), so $c$-stability is equivalent to the statement that
$\alpha$-almost Lewis weights are always within $\alpha^c$ of the true Lewis weights.  Note
that this implies that $(1+O(1/c))$-almost Lewis weights are within $O(1)$ of the true Lewis
weights.
\end{proof}

Now, we prove all the results from Section~\ref{sec:properties}.
\begin{proof}[Proof of Lemma~\ref{lem:fullstable}]
Here, we prove the $\alpha$-almost weight notion of stability, part~\ref{def:stable1} of Definition~\ref{def:stable}.
First, note that for $\alpha$-almost Lewis weights $\ww$,
\begin{equation*}
\textsc{LewisIterate}(\AA, p, 1, \ww) \approx_{\alpha^{p/2}} \ww.
\end{equation*}

In other words, the first step can only change the weights by a power of $p/2$.
But then since $\textsc{LewisIterate}$ is a contraction mapping by $|p/2 - 1|$,
the $t$\textsuperscript{th} step after this, when iterating, can change by only $\alpha^{p/2 |p/2 - 1|^t}$.
The iteration will converge to $\wwbar$, while changing by a total of at most
\begin{equation*}
\alpha^{p/2 \sum_{t=0}^\infty |p/2 - 1|^t} = \alpha^{\frac{p/2}{1 - |p/2 - 1|}}.
\end{equation*}
\end{proof}

\begin{proof}[Proof of Lemma~\ref{lem:sqrtstable}]
Here, we will find it more convenient to use the definition of stability in terms of a reweighted matrix,
part~\ref{def:stable2} of Definition~\ref{def:stable}.

We bound the stability by bounding it infinitesimally.  We let $\rr_i$ and $\ss_i$ be factors that can be
multiplied by the weights of the row, constrained so that $\SS \AA$, when multiplied by
$\WWbar^{1/2-1/p}$ for its Lewis weights $\WWbar$, gives $\RR \AA$.  We will bound the derivative
of $\rr$ with respect to $\ss$.  For any direction vector $\Delta(\ss)$, let the derivative
in that direction by $\Delta(\rr)$.  Then the stability claim is that the max of $\frac{\Delta(\ss)_i}{\ss_i}$
is at most $c$ times the max of $\frac{\Delta(\rr)_i}{\rr_i}$.

It is not apparent how to bound this directly.  However, note that saying that the derivative of $\rr$ with
respect to $\ss$ in the direction of $\Delta(\ss)$ equals $\Delta(\rr)$ is equivalent to saying that the derivative of $\ss$ with
respect to $\rr$ in the direction of $\Delta(\rr)$ equals $\Delta(\ss)$.  Thus, what we want is to show that
that the ratio can't \emph{decrease} too much when differentiating with respect to $\rr$.

For convenience, we will assume that $\AA^T \RR^2 \AA = \II$ (by a change of basis)
and $\norm{\aa_i}^2 = 1$ (by rescaling the rows of $\AA$ simultaneously with $\RR$).

The $\ss$ in terms of the $\rr$ are
\begin{equation*}
\ss_i = \rr_i^{2/p} (\aa_i^T (\AA^T \RR^2 \AA)^{-1} \aa_i)^{1/p - 1/2}.
\end{equation*}

The initial $\ss_i$ are just $\rr_i^{2/p}$.

Then given that $\norm{\aa_i}^2 = 1$ and $\AA^T \AA = \II$, the partial derivative of $\ss_i$ with respect to $\rr_j$ is
\begin{equation*}
\rr_i^{2/p-1} ( (1-2/p) \rr_i \rr_j (\aa_i^T \aa_j)^2 + 2/p \delta_{i,j} ).
\end{equation*}

We may define the symmetric matrix $\MM_{i,j} = ( (1-2/p) \rr_i \rr_j (\aa_i^T \aa_j)^2 + 2/p \delta_{i,j} )$.  Then the partial derivatives gives us
\begin{equation*}
\Delta(\ss) = \RR^{2/p-1} \MM \Delta(\rr).
\end{equation*}

The matrix of $\rr_i \rr_j \rr_j (\aa_i^T \aa_j)^2$ is positive semidefinite--we can write it as $\CC^T \CC$,
where the rows of $\CC$ are scaled tensor squared rows of $\AA$: $\cc_i = \rr_i (\aa_i \otimes \aa_i)$.
Thus, $\MM \succeq 2/p \II$, or equivalently, all eigenvalues of $\MM$ are at least $\frac{2}{p}$.
Thus,
\begin{equation*}
\norm{\MM \Delta(\rr)}_2 \geq \frac{2}{p} \norm{\Delta(\rr)}_2.
\end{equation*}

In other words,
\begin{equation*}
\norm{\RR^{1-2/p} \Delta(\ss)}_2 \geq \frac{2}{p} \norm{\Delta(\rr)}_2.
\end{equation*}

The elements of $\RR^{1-2/p} \Delta(\ss)$ are $\rr_i \frac{\Delta(\ss)_i}{\rr_i^{2/p}}$
$\rr_i \frac{\Delta(\ss)_i}{\ss_i}$.  Finally, since $\sum_i \rr_i^2 = d$,
\begin{equation*}
\norm{\RR^{1-2/p} \Delta(\ss)}_2 \leq \sqrt{d} \max_i \frac{\Delta(\ss)_i}{\ss_i},
\end{equation*}
so that
\begin{equation}
\label{eqn:bigl2}
\max_i \frac{\Delta(\ss)_i}{\ss_i} \geq \frac{2}{p \sqrt{d}} \norm{\Delta(\rr)}_2.
\end{equation}

However, we still need to address the possibility that $\norm{\Delta(\rr)}_2$ is much smaller
than $\max_i \frac{\Delta(\rr)_i}{\rr_i}$.  To do this, consider the $i$ that maximizes that
ratio (and assume by scaling that the ratio is 1).  Then, for that particular $i$,
\begin{equation*}
\frac{\Delta(\ss)_i}{\ss_i} = \frac{2}{p} - (1-2/p) \sum_j \rr_j^2 (\aa_i^T \aa_j)^2 \left ( \frac{\Delta(\rr)_j}{\rr_j} \right ).
\end{equation*}
Note that we can also express $\norm{\Delta(\rr)}_2^2$ as
\begin{equation*}
\sum_j \rr_j^2 \left ( \frac{\Delta(\rr)_j}{\rr_j} \right )^2.
\end{equation*}
Now, $0 \leq \rr_j^2 (\aa_i^T \aa_j)^2 \leq \rr_j^2$, while $\sum_j \rr_j^2 (\aa_i^T \aa_j)^2 = 1$
(since $\AA^T \RR^2 \AA = \II$), so
$\sum_j \rr_j^2 (\aa_i^T \aa_j)^2 \left ( \frac{\Delta(\rr)_j}{\rr_j} \right )$ is at most the average
of $\left ( \frac{\Delta(\rr)_j}{\rr_j} \right )$ in the largest mass-1 set of $j$s.  Thus
it is at most $\norm{\Delta(\rr)}_2$ (since weighted $\ell_2$ norms dominate $\ell_1$ norms when the total
mass sums to 1).
Plugging that in we get
\begin{equation}
\label{eqn:smalll2}
\max_i \frac{\Delta(\ss)_i}{\ss_i} \geq 2/p - (1-2/p) \norm{\Delta(\rr)}_2.
\end{equation}

With $\max_i \frac{\Delta(\rr)_i}{\rr_i} = 1$, we have both
Equation~\ref{eqn:bigl2} and Equation~\ref{eqn:smalll2}.
The former is increasing in $\norm{\Delta(\rr)}_2$ and the latter decreasing,
so the worst case for the minimum of the two of these is when they are equal,
which occurs at $\norm{\Delta(\rr)}_2 = \frac{\frac{2}{p}}{\frac{2}{p \sqrt{d}} + (1-2/p)}$.
This bound is
\begin{align*}
\max_i \frac{\Delta(\ss)_i}{\ss_i} &\geq \frac{\frac{4}{p^2 \sqrt{d}}}{\frac{2}{p \sqrt{d}} + (1-2/p)} \\
&= \frac{4}{2p + (p^2-2p) \sqrt{d}} \\
&= \Omega \left ( \frac{1}{\sqrt{d}} \right ).
\end{align*}
\end{proof}

\begin{proof}[Proof of Lemma~\ref{lem:fullmonotone}]
Assume $p < 2$ (the case for $p=2$ follows from the fact
that $\ell_2$ Lewis weights are just leverage scores).
Let $r$ be the maximum ratio (over $i$) of $\wwbar'_i$ to $\wwbar_i$, which occurs at some row $i^*$.
We want to show than $r \leq 1$; we will suppose that $r > 1$ and obtain a contradiction.

Now, we have
\begin{equation*}
\AA'^T \WWbar'^{1-2/p} \AA' \succeq r^{1-2/p} \AA^T \WWbar^{1-p/2} \AA,
\end{equation*}
since each row of $\AA$ contributes at least $r^{1-2/p}$ as much to the left hand side as to the right.  But the
$\wwbar'_i$ are equal to $(\aa_i^T (\AA'^T \WWbar'^{1-2/p} \AA')^{-1} \aa_i)^{p/2}$, and
\begin{align*}
(\aa_i^T (\AA'^T \WWbar'^{1-2/p} \AA')^{-1} \aa_i)^{p/2} &\leq r^{-p/2 (1-2/p)} (\aa_i^T (\AA^T \WWbar^{1-2/p} \AA)^{-1} \aa_i)^{p/2} \\
&= r^{1-p/2} \wwbar_i
\end{align*}
Thus, in particular, no weight increases by more than $r^{1-p/2}$.  But with $p < 2$, if $r > 1$,
$r^{1-p/2} < r$.  Thus,
\begin{equation*}
\wwbar'_{i^*} \leq r^{1-p/2} \wwbar_{i^*} < r \wwbar_{i^*}
\end{equation*}
This contradicts the fact that $r = \frac{\wwbar'_{i^*}}{\wwbar_{i^*}}$.
\end{proof}

\begin{proof}[Proof of Lemma~\ref{lem:weakmonotone}]
First, assume without loss of generality that $\AA^T \WWbar^{1 - 2/p} \AA$ is equal to the identity matrix.
Then the claim is that all eigenvalues of $\AA'^T \WWbar'^{1 - 2/p} \AA'$ are at least $d^{2/p - 1}$.

For convenience, define $\PP = \AA'^T \WWbar'^{1 - 2/p} \AA'$, and let $\uu$ be an eigenvector of
$\PP$ of minimum eigenvalue, normalized so that $\norm{u}_2 = 1$.  Let $\lambda$ be the eigenvalue
of $\uu$, $\uu^T \PP \uu$.  Then one may see by looking at an orthogonal eigenbasis that
\begin{equation*}
\PP^{-1} \succeq \lambda^{-1} \uu \uu^T.
\end{equation*}

We further define normalized rows $\vv_i = \ww_i^{-1/p} \aa_i$, so that $\AA^T \WWbar^{1 - 2/p} \AA$,
the identity matrix, may be written as
\begin{equation*}
\sum_i \wwbar_i \vv_i \vv_i^T
\end{equation*}
with the $\vv_i$ each being unit vectors.  $\PP$ is lower bounded by its contributions from the original rows
($\AA^T \WWbar'^{1-2/p} \AA$); expressing the latter in terms of the $\vv$ gives
\begin{align*}
\PP &\succeq \sum_i \wwbar_i \frac{(\wwbar'_i)^{1-2/p}}{\wwbar_i^{1-2/p}} \vv_i \vv_i^T \\
&= \sum_i \wwbar_i (\vv^T \PP^{-1} \vv)^{p/2-1} \vv_i \vv_i^T.
\end{align*}

Since $\PP^{-1} \succeq \lambda^{-1} \uu \uu^T$ this further gives
\begin{equation*}
\PP \succeq \sum_i \ww_i \lambda^{1-p/2} (\uu^T \vv)^{p-2} \vv_i \vv_i^T.
\end{equation*}
What we have done is provide a lower bound for the quadratic form $\PP$ given that it has
a small eigenvalue, using the fact that relative leverage scores cannot have decreased
below the leverage scores of the restriction to the eigenspace.  This is not affected by
any increases in the other eigenvalues,

Plugging $\uu$ into the quadratic form implies
\begin{equation*}
\uu^T \PP \uu \geq \lambda^{1-p/2} \wwbar_i (\uu^T \vv)^p.
\end{equation*}

Now, the original quadratic form was the identity, so $\sum_i \ww_i = d$ and $\sum_i \wwbar_i (\uu^T \vv)^2 = 1$.
By a weighted power-mean inequality, $\sum_i \wwbar_i (\uu^T \vv)^p \geq d^{1-p/2}$.

Plugging in $\lambda = \uu^T \PP \uu$ finally gives
\begin{align*}
\lambda &\geq \lambda^{1-p/2} d^{1-p/2} \\
\lambda^{p/2} &\geq d^{1-p/2} \\
\lambda &\geq d^{2/p - 1}.
\end{align*}
\end{proof}
